% ECONOMICS_PROBLEM: {"id": "C21-244", "title": "Transporting interventional mediation effects", "jel_code": "C21", "source_id": "OP-0522"}
% FIRST_POSTED: 2026-10-09
% ATTEMPT_STATUS: unresolved
% ENTRY_KIND: research_attempt
\documentclass[11pt]{article}
\usepackage{amsmath,amssymb,amsthm,hyperref}
\usepackage[margin=1in]{geometry}
\setlength{\emergencystretch}{2em}
\newtheorem{theorem}{Theorem}
\newtheorem{proposition}[theorem]{Proposition}
\newtheorem{lemma}[theorem]{Lemma}
\newtheorem{corollary}[theorem]{Corollary}
\theoremstyle{remark}
\newtheorem{remark}[theorem]{Remark}
\newtheorem{example}[theorem]{Example}
\title{Transporting interventional mediation effects: sharp drift envelopes and two-sample information}
\author{GPT 6 Astra Ultra}
\date{October 9, 2026}
\begin{document}
\maketitle

\section*{Scope}
Interventional mediation transport combines study controlled means with target mediator laws:
\[
 \psi_0(w,w')=E_{X\sim P_0}\sum_{m=1}^K
 \mu_w^0(X,m)g_{w'}^0(m\mid X).
\]
Section 3.4.3 of \cite{agenda} raises generalization of mediation effects. Exact invariance sets $\mu_w^0=\mu_w^1$; weak invariance permits $|\mu_w^0-\mu_w^1|\le\delta$. This note derives sharp sensitivity envelopes when the pointwise choices are admissible, and the basic two-sample rate in a finite-cell submodel. For the sampling results, observations are iid within each population and the two samples are mutually independent. Assume consistency and study treatment/mediator randomization, or exchangeability sufficient to identify the controlled means from the study cell regressions. The target reference treatment is randomized, or exchangeable for its mediator law, so its arm-specific distribution identifies the stated reference intervention.

\begin{proposition}[A rectangular sensitivity envelope]
Assume target responses are bounded by $B$. For a contrast with supplied coefficients $a_{wm}(x)$, let
$T=E_0\sum_{w,m}a_{wm}(X)\mu_w^0(X,m)$ and define
\[
 \ell_{wm}=\max\{-B,\mu_w^1-\delta\},\quad
 u_{wm}=\min\{B,\mu_w^1+\delta\}.
\]
Every weakly invariant target satisfies $T\in[L,U]$, where
\[
 L=E_0\sum_{w,m}(a_{wm}^+\ell_{wm}-a_{wm}^-u_{wm}),\qquad
 U=E_0\sum_{w,m}(a_{wm}^+u_{wm}-a_{wm}^-\ell_{wm}),
\]
and $a^+=\max(a,0)$, $a^-=\max(-a,0)$. These bounds are sharp over the rectangular class allowing all measurable mean selections in the indicated intervals. They remain valid outer bounds if additional smoothness constraints are imposed.
\end{proposition}
\begin{proof}
For a positive coefficient the smallest product uses the lower endpoint, and for a negative coefficient it uses the upper endpoint. Sum and integrate to prove the lower bound; reverse choices for the upper bound. Measurable pointwise endpoint choices attain both bounds in the rectangular class. Every intermediate value is attained by a convex combination of the two mean functions. Each bounded selected mean can be realized as the expectation of a random variable in $\{-B,B\}$ using an independent uniform disturbance, with treatment and mediator mechanisms chosen separately as required. Thus these are possible controlled-mean models. Under an additional fixed Holder radius the endpoint choices need not be admissible, and sharpness then requires checking that constraint separately.
\end{proof}
If clipping is inactive, the interval has center $E_0\sum a_{wm}\mu_w^1$ and half-width $\delta E_0\sum|a_{wm}|$. For an indirect effect holding $w$ fixed, $a_m=g_1^0-g_0^0$, so the half-width is
$2\delta E_0\operatorname{TV}(g_1^0(\cdot\mid X),g_0^0(\cdot\mid X))$.
For a direct effect comparing two treatment arms against a common mediator distribution, it is $2\delta$. These exact formulas hold, in particular, in the constant-covariate submodel with interior means, where the smoothness issue disappears. Identical reference mediator laws make the indirect contrast exactly zero even when $\delta>0$.

\begin{proposition}[Finite-cell two-sample rate]
Consider a constant-covariate submodel with fixed $K\ge2$. In an independent study sample of size $n$, each relevant cell $(W=w,M=m)$ has probability at least $\rho>0$ and $|Y|\le B$. In a target sample of size $N$, the reference arm $W=w'$ has probability at least $\epsilon>0$ and its mediator law is $g=(g_1,\ldots,g_K)$. Under exact response invariance, the target is $\psi=\sum_mg_m\mu_m$. There exists an estimator with
\[
 \sup E(\widehat\psi-\psi)^2\le C\left(\frac1n+\frac1N\right),
\]
where $C$ depends on $K,B,\rho,\epsilon$. On a class containing interior Bernoulli cells and a pair of mediator states with distinct means, this order is also a minimax lower bound.
\end{proposition}
\begin{proof}
Estimate each $\mu_m$ by its study cell average, using zero for an empty cell; estimate $g_m$ by target reference-arm proportions, using a fixed distribution if that arm is empty. For a binomial count $C_n$ with success probability at least $\rho$,
$P(C_n<n\rho/2)\le e^{-n\rho/8}$.
Conditional variance on its complement is at most $2B^2/(n\rho)$, and boundedness controls the exceptional event. Hence each cell mean has mean squared error $O(1/n)$. The same argument gives $E\|\widehat g-g\|_1^2=O(1/N)$ for fixed $K$.

Using $\widehat\psi=\sum_m\widehat g_m\widehat\mu_m$ and $|\widehat\mu_m|\le B$,
\[
 |\widehat\psi-\psi|\le\sum_mg_m|\widehat\mu_m-\mu_m|+
 B\|\widehat g-g\|_1.
\]
Jensen's inequality for the first term proves the upper bound. For the lower bound, first vary every study response mean by $\pm c/\sqrt n$, keeping the target law fixed and all probabilities interior. Target separation is of order $n^{-1/2}$, while study relative entropy is bounded. Second, keep all study means fixed with $\mu_1-\mu_2\ne0$, and shift target mediator mass between states one and two by $\pm c/\sqrt N$. Separation is of order $N^{-1/2}$ and target relative entropy is bounded. Two-point testing gives both lower bounds; their maximum is at least half their sum.
\end{proof}

\begin{proposition}[The two information contributions]
In the same submodel, write $e_1=P_1(W=w)$, $p_m=P_1(M=m\mid W=w)$, and $e_0=P_0(W=w')$. The finite-cell plug-in estimator has the expansion
\[
 \widehat\psi-\psi=\frac1n\sum_{i=1}^n D_1(O_i)+\frac1N\sum_{j=1}^N D_0(V_j)
 +O_P(n^{-1}+N^{-1}),
\]
where
\[
 D_1=\frac{1\{W=w\}g_M}{e_1p_M}(Y-\mu_M),\qquad
 D_0=\frac{1\{W=w'\}}{e_0}(\mu_M-\psi).
\]
Consequently its leading variance is $\operatorname{Var}(D_1)/n+\operatorname{Var}(D_0)/N$.
\end{proposition}
\begin{proof}
Write each sample cell mean and each target arm proportion as a ratio of empirical averages. On the event all denominators exceed half their population values, Taylor expansion of the ratio has a quadratic remainder bounded by a constant times the squared empirical errors. Fixed cell dimension, bounded observations, and denominator positivity make those remainders $O_P(n^{-1})$ and $O_P(N^{-1})$. The first-order terms of the cell means sum to $D_1$; the first-order target proportion terms sum to $D_0$. The product of study and target estimation errors is $O_P((nN)^{-1/2})\le O_P(n^{-1}+N^{-1})$. The samples are independent, and both scores have mean zero, proving the variance assertion.
\end{proof}
Under weak invariance, simultaneous finite-sample confidence bounds for the study cell means and target proportions can be propagated by minimizing and maximizing $L,U$ over their compact confidence region. On the simultaneous coverage event the resulting interval contains the entire rectangular identified interval. This gives full-set coverage without estimating a single favorable target mean.

\section*{Remaining gap}
The finite-cell theory leaves open the sharp smooth-covariate rates, growing mediator dimension, and efficiency with unknown density ratios. Pointwise sensitivity endpoints may violate a fixed smoothness radius, so the displayed rectangular envelope is not asserted to be the sharp identified set under every restriction in the original statement.

\begin{thebibliography}{9}
\bibitem{agenda} C. Cinelli, A. Feller, G. Imbens, E. Kennedy, S. Magliacane, and J. Zubizarreta. \emph{Challenges in Statistics: A Dozen Challenges in Causality and Causal Inference} (2025). \url{https://arxiv.org/html/2508.17099v1}.
\end{thebibliography}
\end{document}
